\section{Threats to Validity}
The reference labels were generated through two AI annotators, a third AI adjudicator,
and a model-visible-context audit, then subsequently confirmed through a confidential
blind human review conducted outside the public repository. The review covered all 413
samples and produced zero disagreements and zero label changes. However, the public
evidence is an operator-confirmed attestation: reviewer identities, reviewer count,
worksheets, and raw materials are withheld for confidentiality and cannot be independently
audited from the repository. Human review occurred after the frozen model evaluation;
the unchanged label SHA preserves the applicability of the existing scoring but does not
make the labels infallible or universally generalizable ground truth.

The authorized challenge set is not natural production traffic and contains no confirmed
successful exploit executions. The primary model's safe FPR was 43/209 (20.57\%) and its
unknown recall was 0, so the model is not suitable for autonomous production blocking.
An observed low or zero rate is a finite-sample result whose residual uncertainty is
quantified by exact binomial intervals, not a population guarantee. Cross-stack latency
(Transformers GPU generation versus Ollama HTTP round-trip) is not directly comparable.
Family Macro-F1 uses reference-attack-family labels only. All eight locked primary systems
completed. Model size is reported only when derived from filesystem stat and tied to an
artifact SHA-256. The merged BF16 safetensors and f16 GGUF share a checkpoint lineage but
are not bit-identical artifacts.
